%% notes-tex/perturbation-operator/sections/1-operators.tex
%% SOURCE: paper/nature-biotech/sections/methods.tex (Eq. ptoken, pert, pertblock);
%% torchcell/models/equivariant_cell_graph_transformer.py on branch
%% multimodal-phenotype-retrospective (class EquivariantPerturbationTransform, the null-sink
%% block and its design comment); notes-tex/figure-3-gate/review/2026-10-04/
%% 05_architecture_and_operator.md for the eval-mode probe.

\section{Three operators that differ in one step\secstatus{todo}}
\label{sec:operators}

\subsection{Symbols\secstatus{todo}}

\Cref{tab:symbols} defines every symbol used in this document. Notation follows the
manuscript's Methods.

\begin{table}[h]
\footnotesize
\caption[Symbols]{\textbf{Symbols, in alphabetical order.} Widths are those of the
1{,}084{,}619-parameter expression model (\file{conf/cgt_expr_012.yaml}); the 6.68 million
parameter model of rounds v13 to v20 has the same operator at the same width.
\src{paper/nature-biotech/sections/methods.tex;
torchcell/models/equivariant_cell_graph_transformer.py}}
\label{tab:symbols}
\begin{tabular}{@{}lp{0.80\linewidth}@{}}
\toprule
symbol & meaning \\
\midrule
$a\in[h]$ & index of an attention head; the operator has $h=9$ \\
$b_K,\,b_V$ & bias vectors of the key and value projections \\
$b_\varnothing$ & the null sink's one learned scalar (\file{null_bias}), shared by all heads \\
$\beta_{i,t}$ & weight gene $i$ puts on perturbation $t$, one scalar per head \\
$d$ & width of a gene state, 90 \\
$d_k=d/h$ & width of one head, 10 \\
$\delta_{i,t}$ & deleted-gene flag: 1 if gene $i$ is the gene that perturbation $t$ deletes, else 0 \\
$\Delta h_i$ & vector the operator adds to gene $i$'s state \\
$g(t)$ & index of the gene that perturbation $t$ deletes \\
$h_i\in\mathbb{R}^d$ & wild-type state of gene $i$, output of the encoder $F_\theta$; the same for every strain \\
$h_i^{\mathrm{pert}}$ & state of gene $i$ in the perturbed strain, output of the operator \\
$k_t=\mathbf{p}_tW_K$ & key of perturbation $t$ \\
$M$ & number of perturbations in a strain; 1 for every expression, proteome and morphology strain used here \\
$N$ & number of genes, 6{,}607 \\
$\mathbf{p}_t\in\mathbb{R}^d$ & token of perturbation $t$. Methods: $\mathbf{p}_t=r(e_t)+\mathbf{t}_{\tau_t}+m_t\mathbf{w}$. Code: $\mathbf{p}_t=h_{g(t)}$; the type embedding $\mathbf{t}_\tau$ and the magnitude term $m_t\mathbf{w}$ are not implemented \\
$q_i=h_iW_Q$ & query of gene $i$ \\
$s_{i,t}=q_i\cdot k_t/\sqrt{d_k}$ & score: relevance of perturbation $t$ to gene $i$ \\
$\sigma(x)=1/(1+e^{-x})$ & logistic sigmoid \\
$v_t=\mathbf{p}_tW_V$ & value of perturbation $t$, the content that can be delivered \\
$W_O$ & output matrix mixing the concatenated heads \\
$W_Q,\,W_K,\,W_V$ & query, key and value projection matrices \\
\bottomrule
\end{tabular}
\end{table}

All three operators share the score, the value, and the residual block that follows,
\begin{equation}
\Delta h_i=\sum_{t=1}^{M}\beta_{i,t}\,v_t,\qquad
\hat h_i=\operatorname{LN}(h_i+\Delta h_i),\qquad
h_i^{\mathrm{pert}}=\operatorname{LN}\!\big(\hat h_i+\operatorname{FFN}(\hat h_i)\big),
\label{eq:block}
\end{equation}
where $\operatorname{LN}$ is layer normalization and $\operatorname{FFN}$ a position-wise
feed-forward network. They differ only in how the score becomes the weight $\beta_{i,t}$.

\subsection{Operator A: softmax over perturbations\secstatus{todo}}

The Methods and the code's default.
\begin{equation}
\beta^{A}_{i,t}=\frac{e^{s_{i,t}}}{\sum_{u=1}^{M}e^{s_{i,u}}},\qquad
\sum_{t=1}^{M}\beta^{A}_{i,t}=1 .
\label{eq:softmax}
\end{equation}
At $M=1$ the single weight is $e^{s_{i,1}}/e^{s_{i,1}}=1$ for every gene, whatever the
score, so $\Delta h_i=v_1$: one vector shared by all 6{,}607 genes. The score cancels, and
the 16{,}200 weights of $W_Q$ and $W_K$ receive zero gradient on single-deletion data.
Measured in evaluation mode, the largest difference between any gene's $\Delta h_i$ and
gene 0's is $0.0$.
\src{notes-tex/figure-3-gate/review/2026-10-04/05_architecture_and_operator.md}

Every expression, proteome and morphology run used operator A except the arms of
\cref{sec:measured}. The gene-interaction models of experiments 025 and 030 use it too,
at $M=2$ or $3$, where the softmax is not degenerate.

\subsection{Operator B: the per-pair sigmoid of Fig.~1f\secstatus{todo}}

\begin{equation}
\beta^{B}_{i,t}=\sigma(s_{i,t})
=\sigma\!\Big(\frac{(h_iW_Q)\cdot(\mathbf{p}_tW_K)}{\sqrt{d_k}}\Big).
\label{eq:sigmoid}
\end{equation}
No sum over $t$ appears. Each weight is set by its own score and lies in $(0,1)$. At $M=1$
the weight differs from gene to gene and $W_Q$, $W_K$ receive gradient. At $M=3$ the three
weights can all be near 1, so effects add, or all near 0, so the gene ignores all three.
Operator B exactly as drawn has not been implemented; the null sink is its closest
implemented relative.

\subsection{Operator C: the null sink\secstatus{todo}}

The null sink (\file{null_sink=true}) prepends one extra key and value to the perturbation
set: an all-zero token, index $\varnothing$, with a learned scalar $b_\varnothing$ added to
its logit. The softmax then runs over $M+1$ entries. The zero token passes through
projections that carry biases, so its key is $b_K$ and its value is $b_V$; the deleted
gene's key is $\mathbf{p}_tW_K+b_K$, and $b_K$ cancels between the two logits. At $M=1$,
\begin{equation}
\beta^{C}_{i,1}=\frac{e^{s_{i,1}}}{e^{s_{i,1}}+e^{b_\varnothing}}
=\sigma\big(s_{i,1}-b_\varnothing\big),\qquad
\Delta h_i=\beta^{C}_{i,1}\,v_1+\big(1-\beta^{C}_{i,1}\big)\,b_V .
\label{eq:sink}
\end{equation}
A two-way softmax is the sigmoid of the difference of its two logits, which is why the null
sink has the form of \cref{eq:sigmoid} at one deletion. \Cref{tab:b-vs-c} lists where the
two part.

\begin{table}[h]
\footnotesize
\caption[Fig.~1f against the null sink]{\textbf{The null sink equals the operator of
Fig.~1f only at one deletion.} For multi-gene strains and for natural isolates it keeps the
softmax's shared budget among the real perturbations.
\src{torchcell/models/equivariant_cell_graph_transformer.py, null-sink design comment}}
\label{tab:b-vs-c}
\begin{tabular}{@{}p{0.20\linewidth}p{0.30\linewidth}p{0.42\linewidth}@{}}
\toprule
 & B (Fig.~1f) & C (null sink) \\
\midrule
threshold & none; $\beta=0.5$ at score 0 & $b_\varnothing$, learned, one scalar for all nine heads \\
what the unused share delivers & nothing & $b_V$, a learned constant vector \\
at $M\ge 2$ & weights independent; their sum lies between 0 and $M$ & a softmax over $M+1$ entries; the weights on real perturbations sum to at most 1, so effects are averaged and never added \\
magnitude & as is & optional rescale by $1/\sigma(-b_\varnothing^{\mathrm{init}})$ (\file{null_sink_magnitude_match}), which restores the reference's norm at initialization \\
\bottomrule
\end{tabular}
\end{table}

Initialization decides whether the sink takes part. With $b_\varnothing=-4$ the gate is
$\sigma(s+4)\approx 0.98$: the sink holds about 2 percent of the mass and the operator is
nearly operator A. In the first attempt the bias moved from $-4.0$ to between $-3.92$ and
$-3.95$ by epoch 187. With $b_\varnothing=0$ the gate starts at $0.5$, and the model has to
close the sink actively for it to be useless.
\src{notes/experiments.019-simb-multimodal.scripts.gh_expr_008_arm.md}

\subsection{Which normalization makes the weights uniform\secstatus{todo}}

Two normalizations are in play and they behave differently.

\pfstep{Normalizing the weights over perturbations restores the degeneracy} If the sigmoid
weights are divided by their sum over $t$,
$\tilde\beta_{i,t}=\beta^{B}_{i,t}/\sum_u\beta^{B}_{i,u}$, then at $M=1$
$\tilde\beta_{i,1}=1$ again. Any normalization across the perturbation axis reproduces
operator A at one deletion. Fig.~1f as drawn has none.

\pfstep{The layer normalization that follows does not} $\operatorname{LN}$ acts on each
gene's vector separately, across its $d$ coordinates, and does not compare genes. Two genes
with $\beta_i=0.1$ and $\beta_j=0.9$ receive different mixtures of their own state and
$v$, with different directions, and $\operatorname{LN}$ preserves direction. It removes
overall magnitude: if $\beta_iv$ is far larger than $h_i$ for every gene, every $\hat h_i$
approaches the direction of $v$ and the genes become alike. That is a matter of scale,
$\lVert v\rVert$ against $\lVert h_i\rVert$, and not a degeneracy.
